% DOTE2011G · In-Class Exercise 2
% Halloween · Jigsaw latch game · one-sided CI
%
% Recommended: xelatex Ex2.tex

\documentclass[11pt,a4paper,UTF8]{ctexart}

\usepackage[margin=2.2cm]{geometry}
\usepackage{amsmath,amssymb}
\usepackage{enumitem}
\usepackage{xcolor}
\usepackage{tcolorbox}
\tcbuselibrary{skins}
\usepackage{fancyhdr}
\usepackage{microtype}

\definecolor{blood}{HTML}{7A0C0C}
\definecolor{night}{HTML}{120C14}
\definecolor{paper}{HTML}{F4EFE6}
\definecolor{amber}{HTML}{C9A227}

\pagestyle{fancy}
\fancyhf{}
\fancyhead[L]{\small DOTE2011G · Statistical Analysis for Business Decisions}
\fancyhead[R]{\small In-Class Exercise 2}
\fancyfoot[C]{\thepage}
\renewcommand{\headrulewidth}{0.4pt}

\begin{document}

\begin{center}
  {\Large\bfseries\color{blood} I Want to Play a Game}\\[0.25em]
  {\large\bfseries Halloween · One-sided confidence interval}\\[0.55em]
  {\normalsize A lower bound for mean latch strength}\\[0.7em]
  {\small Participation exercise · please show your working}
\end{center}

\vspace{0.35em}

\begin{tcolorbox}[
  enhanced,
  colback=paper,
  colframe=blood,
  boxrule=1.1pt,
  arc=2mm,
  left=10pt,right=10pt,top=8pt,bottom=8pt,
  title={\bfseries Jigsaw's message},
  fonttitle=\color{white}\bfseries,
  coltitle=white,
  colbacktitle=blood
]
\textit{``You advertised that the latches in your company have mean breaking strength at least 50.
Many people have bought your products because of your claim.
I brought one, which was WEAK!
Now, let's play a game.
Go pick up 10 latches on the shelf behind you --- they are all your products.
Then you can measure their strengths using the machine in front of you.
If you can show that your claim is true with a 95\% confidence level, I will forgive you for selling me a weak latch and let you go.
Otherwise, you are a liar --- and now you have to pay for it!''}
\end{tcolorbox}

\vspace{0.75em}

\noindent\textbf{Your data.}
You randomly pick 10 latches and measure breaking strengths (kg):
\[
51,\ 55,\ 48,\ 53,\ 57,\ 50,\ 54,\ 49,\ 56,\ 52.
\]
The population standard deviation is \emph{unknown}. Treat the strengths as a random sample from a normal population with mean $\mu$.

\vspace{0.65em}

\noindent\textbf{Task.}
Give a \textbf{95\% lower confidence bound} for $\mu$ (write the interval in the form $[L,\ \infty)$).
Then decide: does this interval support the advertising claim $\mu\ge 50$ at the 95\% level?

\vspace{0.85em}
\noindent\rule{\textwidth}{0.4pt}\\[0.55em]
{\small
No written submission --- a few students will be invited to present in class.
Happy Halloween.
}

\end{document}
